%Plantilla basada en "Template for Masters / Doctoral Thesis" (plantilla disponible en writeLaTex) que subió LaTeXTemplates.com

\documentclass[11pt]{book}
\usepackage[paperwidth=17cm, paperheight=22.5cm, bottom=2.5cm, right=2.5cm]{geometry}
\usepackage{amssymb,amsmath,amsthm} %paquete para símbolo matemáticos
\usepackage{bbm}
\usepackage{physics}
\usepackage[spanish]{babel}
\usepackage[utf8]{inputenc} %Paquete para escribir acentos y otros símbolos directamente
\usepackage{enumerate}
\usepackage{graphicx}
%\usepackage{subfig} %para poner subfiguras
\graphicspath{{Img/}} %En qué carpeta están las imágenes
\usepackage[nottoc]{tocbibind}
\usepackage[pdftex,
            pdfauthor={Germán Eduardo Osorio Leiva},
            pdftitle={Revisión de los fundamentos termodinámicos desde la perspectiva de la teoría cuántica de la información},
            pdfsubject={Física},
            pdfkeywords={entrelazamiento, mecánica estadística},
            pdfproducer={Latex con hyperref},
            pdfcreator={pdflatex}]{hyperref}



\begin{document}

%----------------------------------------------------------------------------------------
%	COMANDOS PERSONALIZADOS
%----------------------------------------------------------------------------------------

%SI TU TESIS TIENE TEOREMAS Y DEMOSTRACIONES, PUEDES DESCOMENTAR Y USAR LOS SIGUIENTES COMANDOS

%\renewcommand{\proofname}{Demostración}
%\providecommand{\norm}[1]{\lVert#1\rVert} %Provee el comando para producir una norma.
%\providecommand{\innp}[1]{\langle#1\rangle} 
%\newcommand{\seno}{\mathrm{sen}}
%\newcommand{\diff}{\mathrm{d}}

%\newtheorem{teo}{Teorema}[section] 
%\newtheorem{cor}[teo]{Corolario}
%\newtheorem{lem}[teo]{Lema}

%\theoremstyle{definition}
%\newtheorem{dfn}[teo]{Definición}

%\theoremstyle{remark}
%\newtheorem{obs}[teo]{Observación}

%\allowdisplaybreaks


%----------------------------------------------------------------------------------------
%	PORTADA
%----------------------------------------------------------------------------------------

\title{Revisión de los fundamentos termodinámicos desde la perspectiva de la teoría cuántica de la información} %Con este nombre se guardará el proyecto en writeLaTex

\begin{titlepage}
\begin{center}

\textsc{\Large Universidad de los Andes}\\[4em]

%Figura
\begin{figure}[h]
\begin{center}
\includegraphics[scale=0.045]{ellogo.png}
\end{center}
\end{figure}

\vspace{4em}

\textsc{\huge \textbf{Revisión de los fundamentos termodinámicos desde la perspectiva de la teoría de la información cuántica}}\\[4em]

\textsc{\large Tesis}\\[1em]

\textsc{que para obtener el título de}\\[1em]

\textsc{FÍSICO}\\[1em]

\textsc{presenta}\\[1em]

\textsc{\Large 	GERMÁN EDUARDO OSORIO LEIVA}\\[1em]

\textsc{\large Asesor: }

\end{center}

\vspace*{\fill}
\textsc{Bogotá, D.C. \hspace*{\fill} 2017}

\end{titlepage}


%----------------------------------------------------------------------------------------
%	DECLARACIÓN
%----------------------------------------------------------------------------------------

\thispagestyle{empty}
\vspace*{\fill}
\begingroup
``Con fundamento en los artículos 21 y 27 de la Ley Federal del Derecho de Autor y como titular de los derechos moral y patrimonial de la obra titulada ``\textbf{TÍTULO DE LA TESIS}'', otorgo de manera gratuita y permanente al Instituto Tecnológico Autónomo de México y a la Biblioteca Raúl Bailléres Jr., la autorización para que fijen la obra en cualquier medio, incluido el electrónico, y la divulguen entre sus usuarios, profesores, estudiantes o terceras personas, sin que pueda percibir por tal divulgación una contraprestación''.

\centering

\hspace{3em}

\textsc{AUTOR}

\vspace{5em}

\rule[1em]{20em}{0.5pt} % Línea para la fecha

\textsc{Fecha}
 
\vspace{8em}

\rule[1em]{20em}{0.5pt} % Línea para la firma

\textsc{Firma}

\endgroup
\vspace*{\fill}


%----------------------------------------------------------------------------------------
%	DEDICATORIA
%----------------------------------------------------------------------------------------

\pagestyle{empty}
\frontmatter


%----------------------------------------------------------------------------------------
%	AGRADECIMIENTOS
%----------------------------------------------------------------------------------------

\chapter*{Agradecimientos}
%\markboth{AGRADECIMIENTOS23}{AGRADECIMIENTOS} % encabezado 
Gracias a mis padres.



%----------------------------------------------------------------------------------------
%	PREFACIO
%----------------------------------------------------------------------------------------

\chapter*{Prefacio}

\pagestyle{plain}
%\markboth{PREFACIO23}{PREFACIO} % encabezado 

PUEDEN QUITAR ESTA PARTE


%----------------------------------------------------------------------------------------
%	TABLA DE CONTENIDOS
%---------------------------------------------------------------------------------------

\tableofcontents


%----------------------------------------------------------------------------------------
%	TESIS
%----------------------------------------------------------------------------------------
\mainmatter %empieza la numeración de las páginas
\pagestyle{headings}

%  Incluye los capítulos en el folder de capítulos


\chapter{El entrelazamiento y la mecánica estadística}

La idea principal que se quiere mostrar es cómo se puede reemplazar el postulado de probabilidades iguales por un principio canónico general basado en el entrelazamiento cuántico, basado en el entrelazamiento del sistema y el ambiente.

La mecánica clásica nos habla de un sistema físico definido que para todos los tiempos está especificado. Este sistema evolucionan de manera determinista. Pero lo que sorprende al tratar sistemas termodinámicos es que aunque se hable de un sistema clásico este puede mostrar propiedades que dependan de promedios estadísticos. La conexión que hay entre el determinismo y estas probabilidades es una discusión que lleva desde los inicios de la termodinámica. Aunque el planteamiento que se mostrará no necesita de los métodos típicos de la mecánica estadística como por ejemplo: aleatoriedad subjetiva, promedio sobre el ensamble o promedio temporal; Esto le da una fuerza a estas ideas ya que no debe entrar en problemas de ergodicidad. 

Este nuevo enfoque se tiene un universo que está compuesto por el sistema y el ambiente. Dando como condición que el ambiente sea los suficientemente grande. Este universo está descrito por un estado cuántico puro ( se conoce el estado de manera exacta) que obedece una restricción global, que el sistema alcance el equilibrio térmico por medio de la interacción mutua (termalización) es un producto del entrelazamiento del sistema y el ambiente (Popescu et al,2006). Más adelante se le dará una definición más rigurosa que ayudará a dar cotas para la expresión "ambiente suficientemente grande". Con este enfoque se quiere formular un principio canónico general  el sistema estará termalizado para casi todos los estados puros del universo, dando límites cuantitativos. La restricción que se impone no es una específica esto generaliza los resultados tradicionales dados en la literatura donde se toma por restricción la energía (cualquier libro, año).
Estos resultados no miran la evolución del sistema pero debido a que la mayoría de los estados del universo están termalizados se preve que no importa el estado inicial cualquier evolución llevará al estado a uno en el equilibrio.

(
Lo que mi asesor no me ha respondido.
)

Ya poniendo las ideas más explícitas se supone tener un sistema cuántio aislado y grande que se llamará el universo, que se descompone en dos partes el sistema $S$ y el ambiente $E$. La dimensión del ambiente es mucho maś grande que la del sistema $S$. Además se le impone una restricción global al universo llamada $R$. Desde la mecánica cuántica estas condiciones se pueden poner como restricciones en el espacio de Hilbert, restricción de los estados posibles:

\begin{equation}
\mathcal{H_{R}}\subseteq \mathcal{H_{S}}\otimes \mathcal{H_{E}},
\end{equation}

Donde $\mathcal{H_{S}}$ y $\mathcal{H_{E}}$ son los espacios de Hilbert del sistema y el ambiente con dimension $d_{\mathcal{S}}$  y $d_{\mathcal{E}}$ respectivamente. Es bueno recalcar que $R$ es una restricción arbitraria generalmente se toma como la energía del universo.
Ahora se define el estado equiprobable del universo bajo $R$ como

\begin{equation}
\mathcal{E_{R}} = \frac{1}{d_{R}} \mathbbm{1_{\mathcal{R}}},
\end{equation}

Donde $\mathbbm{1_{\mathcal{R}}}$ es el operador identidad (proyección) sobre el espacio de Hilbert  $\mathcal{H_{R}}$  que tiene dimensión $d_{\mathcal{R}}$. Esto se relaciona con el principio de probabilidades iguales porque este es el estado máximamente mezclado en $\mathcal{H_{R}}$ por ser así todos los estados bajo la restricción de R tienen la misma probabilidad de salir.

Definimos $\Omega_{\mathcal{S}}$ como el estado canónico que está restringido por $R$ cuando el universo se encuentra en el estado $\mathcal{E_{R}}$.Esto significa que si se hace una traza parcial al universo sobre el ambiente da como resultado el estado canónico:
\begin{equation}
 \Omega_{\mathcal{S}} =Tr_{\mathcal{E}} \mathcal{E_{R}}.
\end{equation}

Aquí se hace un supuesto importante y es que el universo está en un estado puro $\ket{\phi}$ y no en un estado mixto $\mathcal{E_{R}}$, esto quiere decir que se conoce todo lo que es permitido por la mecánica cuántica del universo si estuviese en un estado mixto significaría que nosotros no tenemos toda la información que se pudiese tener. Ahora lo que se quiere ver es que aun que el estado del universo sea puro el estado reducido del sistema,

\begin{equation}
\rho_{S}=Tr_{E}\ket{\phi}\bra{\phi} 
\end{equation}

se acerca al estado canónico para la gran mayoría de los casos 

\begin{equation}
\rho_{S} \approx \Omega_{S}.
\end{equation}

O sea que para casi todos los estados puros $\ket{\phi} \in \mathcal{H_{R}}$ del universo el sistema se comporta como si el universo estuviese en el estado mixto equiprobable $\mathcal{E_{R}}$. Siendo este el principio general canónico.

Clarificando la idea, el estado canónico $\Omega_{S}$ del sistema es el estado del sistema cuando el universo se encuentra en el estado equiprobable $\mathcal{E_{R}}$. Se puede interpretar el principio general canónico como un principio que estipula que las probabilidades iguales del sistema son aparentes. Porque para casi cualquier estado del universo que sea puro, un subsistema de este universo que cumpla con ser lo suficientemente pequeño se comporta como si el universo estuviese en el estado equiprobable $\mathcal{E_{R}}$. Cabe recordar que aún no se ha especificado la restricción $R$ entonces todo este análisis es general, la restricción no necesariamente debe ser la energía u otras cantidades que se conserven. Esto hace que $\Omega_{S}$ no  deba ser obligatoriamente es el estado canónico usual, puede ser el gran canónico o cualquier otro que sea acorde con la restricción. Este principio puede ser de utilidad cuando la interacción entre el ambiente y el sistema no es débil como siempre se toma.

\section{formulación matemática}

Ahora se procederá a especificar lo dicho anteriormente en un contexto matemático. Lo primordial es decir cuál será la distancia que usaremos para darle un sentido de cercanía a los estados $\rho_{S}$ y $\Omega_{S}$. La distancia a usar es una bastante conocida en la teoría cuántica de la información, la distancia de traza. Esta se define como:
\begin{equation}
D(\rho_{S}, \Omega_{S})= \frac{1}{2} Tr |\rho_{S} -\Omega_{S}|=\frac{1}{2}Tr\sqrt{(\rho_{S} -\Omega_{S})^{\dag}(\rho_{S} -\Omega_{S})}.
\end{equation}

Esta medida cuantifica que tan difícil es diferenciar $\rho_{S}$ y $\Omega_{S}$ por mediciones cuánticas. Un poco de notación $\langle . \rangle$ es el promedio sobre todos los estados $\ket{\phi} \in \mathcal{H_{R}}$ de acuerdo a la medida estandar(unitariamente invariante). Esta medida se usa para hallar volúmenes de conjuntos de estados. Luego por esto se sabe que $\Omega_{S} = \langle \rho_{S} \rangle $.

Entonces el teorema central de este capítulo es:

\textbf{Teorema:}Para un estado escogido de manera aleatoria $\ket{\phi} \in \mathcal{H_{R}} \subseteq \mathcal{H_{S}} \otimes \mathcal{H_{E}} $ y un $\epsilon > 0$ arbitrario, la distancia entre la matriz densidad reducida del sistema $\rho_{S}=Tr(\ket{\phi} \bra{\phi})$  y el estado canónico $\Omega_{S}=Tr \mathcal{E_{R}}$ esta dado probabilísticamente por 

\begin{equation}
Prob[  \norm{\rho_{S} -\Omega_{S}}_{1} \geq \eta ] \leq \eta',
\end{equation}

Donde 

\begin{equation}
\norm{\rho_{S}-\Omega_{S}}_{1}=2 D(\rho_{S}, \Omega_{S})
\end{equation}

\begin{equation}
\eta= \epsilon + \sqrt{ \frac{d_{S}}{d_{E}^{eff}} },
\end{equation}

\begin{equation}
\eta'=2exp(-Cd_{R}\epsilon^{2}).
\end{equation}

y las constantes son: $ C=(18\pi^{3})^{-1}, d_{R} = dim \mathcal{H_{R}}, d_{S} = dim \mathcal{H_{S}} $. $d_{E}^{eff}$ es la medida efectiva del tamaño del ambiente,
\begin{equation}
d_{E}^{eff}= \frac{1}{Tr \Omega_{E}^{2}} \ge \frac{d_{R}}{d_{S}}.
\end{equation}

Donde $\Omega_{E}= Tr_{S} \mathcal{E_{R}}$. Ambas cantidades $\eta $ y $\eta'$ serán pequeñas esto implica que el estado estará cercano al estado canónico con alta probabilidad cuando la dimensión efectiva del ambiente sea mucho más grande que la del sistema ( es decir $d_{E}^{eff} >> d_{S}$) y  $d_{R}\epsilon^2>>1>>\epsilon$. Esta última condición se puede asegurar cuando el espacio acesible total es grande ( es decir $d_{R}>>1$), escogiendo $\epsilon=d_{r}^{-\frac{1}{3}}$.\\

Ya con esto se le da un significado cuantitativo de lo que se ha ido explicando hasta ahora. Además también se dan unas cotas explícitas a cuando se habla de la "mayoría" de los estados. Se tiene una cota exponencialmente pequeña  sobre la probabilidad de encontrar un estado lejano del canónico. (Dar comentario sobre lo intuitivo de las cotas)
\section{Lema de Levy}
El teorema anterior tiene como base principal el lema de Levy (referencias). Este dice que al seleccionar un punto $\phi$ aleatoriamente de una hiperesfera de dimensión alta y que $f(\phi)$ no cambie muy rápido, entonces $f(\phi) \approx \langle f \rangle $ con alta probabilidad, Más exactamente:

\textbf{Lema de Levy:} Dada una función $f: \mathbb{S}^d \to \mathbb{R} $ definida en la hiperesfera $\mathbb{S}^d$, y un punto $\phi \in \mathbb{S}^d $ es escogido de manera aleatoria uniforme, 
\begin{equation}
Prob[|f(\phi)- \langle f \rangle| \ge \epsilon ] \le 2exp(-\frac{2C(d+1)\epsilon^2}{\eta^2})
\end{equation}
donde $\eta$ es la constante de Lipschitz de $f$, dado por $\eta= \sup|\nabla f|$ y $C=(18\pi^3)^-1$.







\thispagestyle{empty}
\chapter{Conclusiones}

En este escrito se hizo un recuento histórico de la termodinámica para llegar a la mecánica estadística. Desde esta perspectiva histórica se mostró la evolución de los distintos conceptos que han fundamentado la mecánica estadística. Se expuso la perspectiva de Boltzmann junto con los las dificultades que esta acarrea. Como las objeciones dadas a su función H y la relación con la entropía termodinámica. La siguiente perspectiva que se exhibió fue la de Gibbs. La cual se condensa en la idea del ensamble. También se señaló los inconvenientes  que el concepto de copias imaginarias tiene, hablar de microestados que posiblemente pueden ser infinitos.\\
Junto con esto se muestra que hay otras perspectivas más modernas entre ellas la de Jaynes y su máxima entropía. La cual forma el primer vínculo entre la mecánica estadística y la teoría de la información. La perspectiva de Jaynes mejora respecto a las anteriores porque no introduce hipótesis a priori y especifica que la probabilidad es vista desde un sentido subjetivo, algo que las otras perspectivas no especificaban.
\\
Ya más adelante se siguió el artículo de Popescu et al  \cite{Popescu2006} mostrando la tipicidad canónica con ayuda de las herramientas de la teoría de la información cuántica. Se replicaron los cálculos del artículo para explicar los detalles del formalismo que los autores propusieron. Con esto puesto se vio la forma en que se conciliaban las perspectivas de Gibbs y Boltzmann gracias a este enfoque. Además el enfoque muestra que el principio de probabilidades a priori se puede reemplazar por el principio general canónico. En esto se concluye que es muy poco probable encontrar un estado que no se encuentre en el estado canónico.
\\
Para terminar se vio las extensiones que tiene los nuevos fundamentos gracias a la tipicidad canónica. Se vio la capacidad del principio general canónico tiene para poder explicar cómo se termaliza un sistema. Se reprodujo el análisis hecho por Linden et al \cite{LindenPaper} en este tema. Aunque no pudieron demostrar la independencia del subsistema, abrieron un campo a nuevas investigaciones.
\thispagestyle{empty}


%----------------------------------------------------------------------------------------
%	APÉNDICES
%----------------------------------------------------------------------------------------

\addtocontents{toc}{\vspace{2em}} % Agrega espacios en la toc

\appendix % Los siguientes capítulos son apéndices

%  Incluye los apéndices en el folder de apéndices

\chapter{Álgebra Lineal}

Cosas.
\thispagestyle{empty}
%\include{Apendices/AppendixB}
%\include{Apendices/AppendixC}

\addtocontents{toc}{\vspace{2em}} % Agrega espacio en la toc


%----------------------------------------------------------------------------------------
%	BIBLIOGRAFÍA
%----------------------------------------------------------------------------------------
\backmatter
\nocite{*}
\bibliographystyle{plain}
\bibliography{bibliografía.bib} %Aquí ponen el nombre del archivo .bib






\end{document}